\documentclass[a4paper, 10pt]{article}
\usepackage{palatino}
\usepackage[a4paper]{geometry}
\usepackage[UKenglish]{babel}
\usepackage[utf8]{inputenc}
\usepackage{afterpage}
\usepackage{fancyhdr}
\usepackage{amsmath}
\usepackage{amstext}
\usepackage{amssymb}
\usepackage{physics}
\usepackage{derivative}
\usepackage{longtable}
\usepackage{tabularx}
\usepackage{dcolumn}
\usepackage{graphicx}
\usepackage{enumerate}
\usepackage{multicol}
\usepackage{pdfpages}
\usepackage{lastpage}
\usepackage{siunitx}
\usepackage{listings}
\usepackage{tikz}
\usepackage{xspace}
\usepackage{caption} 

\AtBeginDocument{\RenewCommandCopy\qty\SI}
\sloppy

\geometry{margin=25mm, top=20mm, bottom=20mm}
\captionsetup{format=plain,indention=0mm, font=small, labelfont=bf, labelsep=period}

\newcommand{\logo}{\begin{figure}[!t] \vspace*{-10mm} \hspace*{-10mm}
\includegraphics[height=23mm]{../defs/usn-logo-en.pdf} \end{figure}}

\newcommand{\labfoot}[1]{\fancyhf[FC]{ \textit{TSE2280, Lab exercise  #1, spring 2024} } }

\newcommand{\regneovinglosningfoot}[1]{\fancyhf[FC]{TSE2280 Regneøving #1 - Løsning }}
\newcommand{\regneovingtitle}[1]{\section*{TSE2280, våren 2024. \qquad Regneøving #1 }}
\newcommand{\regneovinglosningtitle}[1]{\section*{TSE2280, Måleteknikk og signalbehandling \\Løsning regneøving #1 }}

\newcommand{\Fig}[1]{Fig.~\ref{fig:#1}}
\newcommand{\tableref}[1]{Table~\ref{tab:#1}}

\newcommand{\ed}[1]{\ensuremath{\cdot10^{#1}}}
\renewcommand{\deg}{\ensuremath{^\circ}}
\newcommand{\omegahat}{\hat \omega}

\newcommand{\zm}[1]{ z^{-{#1} } }

\newcommand{\ew}[1]{e^{{#1}j \omegahat }}
\newcommand{\ewn}[2]{e^{{#1}j{#2}} }
\newcommand{\Hw}[1]{H\left( e^{{#1}j \omegahat} \right)}
\newcommand{\Hwn}[2]{H \left( e^{{#1}j {#2}} \right)}

\newcounter{problem}
\newcommand{\problem}[1]{\stepcounter{problem}\section*{Problem \arabic{problem}. {#1}}} 

\newcommand{\subproblem}{\begin{enumerate}[a)]}
\newcommand{\subproblemend}{\end{enumerate}}

%=== FIGURE DEFINITIONS ============
\newenvironment{stdfig}{\begin{figure}[!b]\begin{center}}{\end{center}\end{figure}}
\newcommand{\stdgraphics}[1]{\begin{center}\includegraphics[width=0.95\textwidth]{#1}\end{center}}
\newcommand{\scaledgraphics}[2]{\begin{center}\includegraphics[width={#1}\textwidth]{#2}\end{center}}
\newcommand{\includescanFB}[1]{\begin{center}\includegraphics[width=0.90\textwidth]{#1}\end{center}}
\newcommand{\includescanMcC}[1]{\includegraphics[width=0.80\textwidth]{#1}}
\newcommand{\includedrawing}[1]{\includegraphics[scale=0.65]{#1}}
\newcommand{\Matlabfigure}[1]{\includegraphics[scale=0.80]{#1}}

\newcommand{\losning}[1]{\includegraphics[width=0.85\textwidth]{./losninger/#1}}
\newcommand{\losningfig}[1]{\includegraphics[width=0.85\textwidth]{./figs/#1}}
\newcommand{\losningfigfull}[1]{\includegraphics[width=0.99\textwidth]{./figs/#1}}
\newcommand{\losningfighalf}[1]{\includegraphics[width=0.49\textwidth]{./figs/#1}}
\newcommand{\losningfigreduced}[1]{\includegraphics[width=0.49\textwidth]{./figs/#1}}

%=== REFERENCING ===========================================
\newcommand{\Eq}[1]{(\ref{eq:#1})}

% Headers and footers
\pagestyle{fancy}
\renewcommand{\headrulewidth}{0pt}
\renewcommand{\footrulewidth}{0pt}

\fancyhf{} 
\fancyhf[FL]{\textit{\small LH, \today} }
\fancyhf[FR]{Page \thepage\ of \pageref{LastPage}} 


% Style for program listing
\definecolor{codegreen}{rgb}{0, 0.5, 0}
\definecolor{codegray}{rgb}{0.5, 0.5, 0.5}
\definecolor{codepurple}{rgb}{0.58, 0, 0.82}
\definecolor{backcolour}{rgb}{0.98, 0.98, 0.95}

\lstdefinestyle{pythonstyle}{
	language=python,	
	belowcaptionskip=1\baselineskip,
	breaklines=true,
	frame=single,
	numbers=none, 
	breaklines=true,				 
	keepspaces=true,		
	breakatwhitespace=false,		 		 	
	basicstyle=\small\ttfamily,
	keywordstyle=\color{blue},
	stringstyle=\color{codepurple},	
	commentstyle=\color{codegreen},
	identifierstyle=\color{black!95!white},
	backgroundcolor=\color{backcolour},
	tabsize=4,	
	}
\labfoot{2}
\hyphenpenalty=5000
\setcounter{secnumdepth}{3}

\begin{document}
\logo

\title{Spectra, spectrograms, and aliasing} 
\author{TSE2280 Signal Processing}
\date{Lab 2, spring 2025}
\maketitle

\thispagestyle{fancy}	

\section{Introduction}
\suppressfloats[t]
\subsection{Background and Motivation}
This lab is based on two labs that that accompany the course text-book \emph{DSP First} by McClellan et~al.~\cite{mcclellan_dsp_2016},  \emph{Lab P-4: AM and FM Sinusoidal Signals} and \emph{Lab S-8: Spectrograms: Harmonic Lines \& Chirp Aliasing}\cite{mcclellan_lab_2016}. 
The lab demonstrates concepts from Chapters 3 and 4 in the text-book, covering spectra, spectrograms, and aliasing.  The lab will also give practical experience with some common signals, such as square and triangle waves, beats, and chirps.  
The original labs have been modified and converted from Matlab to Python. 

Listening to the sound of a signal can be informative. The sound of the signals can be played using the sound card in the computer and compared to what we see in the spectra and spectrograms. 



\subsection{Software Tools: Python with Spyder and \jupyterlab}
Python is used for for programming, with \spyder~\cite{raybaut_spyder_2024} as the programming environment, and \jupyterlab for reporting.
The signals are represented as \numpy arrays and plotted in \matplotlib. 
This lab also introduces spectral analysis tools from the signal processing modules in \scipy, \texttt{scipy.fft} and \texttt{scipy.signal}, and the module \texttt{sounddevice} to play the sounds over the sound card.

The recommended way to import the Python modules is shown in \tableref{tab:import-libraries}. 
You are free to do this in other ways, but the code examples in this text assumes modules are imported as described here.

\afterpage{\clearpage
\input{../figs/table-import-libraries}}


\section{Theory with Programming Examples}

\subsection{Sinusoids}
Lab 1 introduced sinosoidial waveforms represented as 
\begin{align}
	x(t)&= A \cos(\omega t + \phi) = \Re\left\{ A e^{j(\omega t + \phi)} \right\} 
		=  \Re\left\{ X e^{j\omega t} \right\} & &  , &
	X &= A e^{j\phi} 
\end{align}
where $A$ is the amplitude, $\omega = 2\pi f$ is the angular frequency, $f$ is the frequency, and $\phi$ is the phase. $X$ is a \emph{phasor} or \emph{complex amplitude} that includes both the amplitude and the phase of the signal.
In Lab 1, you also wrote a function to create a signal $x_s(t)$as a sum of several sinusoids described by frequencies $f_k$ and complex amplitudes $X_k$
\begin{align}
	x_s(t)&= \sum_{k=1}^{N}  A_k \cos(2\pi f_k t +\phi_k)
	= \Re\left\{ \sum_{k=1}^{N}  X_k  e^{j2\pi f_k t } \right\}  \quad , & 
	X_k&= A_k e^{j\phi_k}    \:.
	\label{eq:harmonicsum}
\end{align}

\subsection{Beat and Amplitude Modulated Signals}
A beat signal is a sum of two sinusoids with frequencies $f_1$ and $f_2$ that are close. Chapter 3 in the textbook \cite{mcclellan_dsp_2016} shows how this can be interpreted as a sinusoid with a high frequency $f_c=\tfrac{1}{2} (f_2+f_1)$ enclosed in a slowly varying envelope with frequency $f_\Delta=\tfrac{1}{2}(f_2-f_1)$, 
\begin{align}
	x_b(t) &= \cos(2\pi f_1 t) + \cos(2\pi f_2 t) = 2\cos(2\pi f_\Delta t)\cos(2\pi f_c t)  \:.
\end{align}	
Amplitude modulated (AM) signals used in radio trasnmission can be viewed as an extension to the beat signal. The AM-signal is written as
\begin{align}
	x_{AM}(t) &= \left(1+M \cos(2\pi f_0 t) \right)\cos(2\pi f_c t) %\nonumber \\ 
	= \cos(2\pi f_0 t) + \tfrac12\cos(2\pi (f_c-f_0) t) + \tfrac12\cos(2\pi (f_c+f_0) t)  \:.
\end{align}	

\subsection{Frequency Modulated Signals, Chirps}
In a sinusoid with constant frequency, the argument of the cosine-function $\Psi = 2\pi f t + \phi$ varies linearly with time. The time derivative of $\Psi$ is its angular frequency $\omega =   \derivative{\Psi}{t} = 2 \pi f$.
This can be generalised by setting the argument to the cosine-function to an arbitrary function $\Psi(t)$,
\begin{align}
	x_{FM}(t) &= A \cos(\Psi(t)) \;,
\end{align}
and defining the instantaneous frequency as 
\begin{align}
	\omega_i &= \derivative{\Psi(t)}{t} \:, & f_i &= \frac{\omega_i}{2\pi} = \frac{1}{2\pi} \derivative{\Psi(t)}{t} \:.
\end{align}
A \emph{linear chirp}, often just called a \emph{chirp}, is a sinusoid where the frequency changes linearly with time. The instantaneous phase $\Psi(t)$ for a linear chirp is a quadratic function of time, 
\begin{align}
	\Psi(t) &= 2 \pi \mu t^2 + 2 \pi f_0 t + \phi  \:,
\end{align}
and the instantaneous frequency is 
\begin{align}
	f_i &=  2\mu t + f_0 \:.
\end{align}
This is an example of a frequency modulated (FM) signal. 
The linear variation of the frequency can produce an audible sound similar to a siren or a chirp, giving this class of its name.

An example of Python code to generatte a chirp is given in \tableref{tab:example-chirp-code}.

\input{../figs/table-example-chirp-code}

\subsection{Spectra and spectrograms}

\subsubsection{Spectrum and Fourier Coefficients}
The spectrum of a signal is a representation of the frequencies present in the signal and can be calculated by finding the Fourier coefficients $a_k$. For a sampled signal $x[n]$, this can be done efficiently by using the FFT-algorithm, which is available as the function \verb|fft| in the \scipy module also called \verb|fft|. 
Some  remarks about how the results returned from the FFT-algorithm are related to the Fourier coefficients $a_k$ are:
\begin{enumerate}[1.]

\item FFT returns coefficients as complex numbers, with magnitude and phase .
	
\item The result returned from FFT are the Fourier coefficients $a_k$ multiplied by the number of samples $n_s$. 
	The coefficients $a_k$ are found by dividing the FFT output by the number of samples. 
	
\item FFT returns both positive and negative frequency components. The number of frequency values returned from FFT is twice the number of samples in the input vector.

\item The Fourier coefficients returned from FFT are arranged with the negative coefficients \emph{after} the positive coefficients. The helper function \verb|fftreshape| reorganize the results so the negative frequencies come first.

\item  The frequencies corresponding to the Fourier coefficients are returned from the function \verb|fftfreq|.

\item The FFT algorithm works by breaking down the number of points in prime factors, it is most efficient if the number of points contain only one prime factor.
This prime factor is often $2$, and FFTs are often evaluated for sequences of $2^n$ points, e.g., 256, 512, 1024, 2048, etc...
The mathematical details can be found in the text-book \cite{mcclellan_dsp_2016}.
	
\end{enumerate}

Example code for how to calculate the Fourier coefficients, order them correctly, and plot the magnitude and phase is shown in \tableref{tab:calculate-fft-code}.

\input{../figs/table-fft-code}


\subsubsection{Power Spectrum}
The power associated with frequency component $k$ is $\tfrac{1}{2}|a_k|^2$. This can be expressed as a \emph{power spectral density, PSD}, or power per frequency interval, by normalising with the frequency interval $\Delta f=1/T_0$. The unit of the PSD is \emph{amplitude square per frequency}, e.g., if $x(t)$ is measured in \unit{Volt}, PSD has unit \unit{V^2/Hz}. This scaling makes the PSD independent of $T_0$ and sample rate $f_s$, making comparison between spectra easier.

For a real-valued signal, the positive and negative frequency coefficients are complex conjugates, $a_{-k}=a_k^*$ and $|a_{-k}|^2=|a_k|^2$. In this case, the power in the negative and positive frequencies can be added to a \emph{single-sided PSD}. 
This is the most common way of presenting power spectra, the power spectral density $P_{xx}$ is expressed by the Fourier coefficients as 
\begin{align}
	P_{xx} = \begin{cases} a_0^2	& k=0 \\ \tfrac{1}{2} |a_k|^2 	& k >0 \end{cases}  \qquad .
\end{align}
This theory is built into a \scipy function \verb|scipy.signal.periodogram|. All all we need to do to find the correctly scaled power spectral density of a signal is to call this function, which also returns the frequency vector. Example code for this is given in \tableref{tab:psd-code}. The calculation of the PSD is only one line, the rest is for plotting the results. 

\input{../figs/table-psd-code.tex}

\subsubsection{Spectrogram}
A spectrum that changes with time can be illustrated with a spectrogram, see Chapter~3-6 in the course text-book \cite{mcclellan_dsp_2016}.
The spectrogram is found dividing the time signal in short intervals of length $T_0$, finding the PSD over each segment, and plotting the result on atwo-dimensional intensity plot with time and frequency on the axes.
Some comments to the use of spectrograms are
\begin{enumerate}
	\item Python offers several versions of spectrogram. 
	One version that is straightforward and simple to use is found in \scipy, \verb|scipy.signal.spectrogram|. 

	\item Spectrograms are numerically calculated over short segments of the signal and give an estimate of the time-varying frequency content of the signal. The finite length of the time segments limits the frequency resolution in the spectrogram, see \chapterref{tf-resolution}.
	
	\item Spectrograms can be difficult to configure.
	Critical parameters are the length of segments and the scales of the frequency axis and intensity display.	
	Different settings can create spectrograms that look different, although presenting the same data.

\end{enumerate}
Example of code for plotting a spectrogram is shown in \tableref{tab:spectrogram-code}.

\input{../figs/table-spectrogram-code.tex}

\subsection{Resolution in Time and Frequency} \label{tf-resolution}
A signal $x(t)$ over an interval $T_0$ can be written as a sum of complex exponentials where the frequencies are the harmonics of the fundamental $f_k= kf_0=k/T_0, k=0, 1, 2, \ldots$, 
\begin{align}
	x(t) &= \sum_{k=-\infty}^{+\infty} a_k e^{j2\pi k t/T_0 }  \:.
	\label{eq:fourier-synthesis}
\end{align}
The strength of the frequency components are given by the Fourier coefficients $a_k$.
They can be found from the signal $x(t)$ by the following equation from the text-book \cite{mcclellan_dsp_2016} 
\begin{align}
	a_k &= \frac{1}{T_0} \int_{0}^{T_0} x(t) e^{-j2\pi k t/T_0 } \di t 	\:.
	\label{eq:fourier-coefficients}
\end{align}
This gives an important relation between resolution in time and frequency. The frequencies of the Fourier coefficients are $0, f_0, 2f_0, 3f_0, \ldots$, so the spacing between them is $\Delta f = f_0 = 1/T_0$. 
A good resolution in frequency, a low $\Delta f$, requires a long observation time, a large $T_0$. The segment length $T_0$ is the resolution in time, $\Delta t = T_0$. 
This gives the important relation between resolution in time and frequency,
\begin{align}
	\Delta f &= 1/T_0 	&	\Delta t &= T_0	&	\Delta f &= \frac{1}{\Delta t} 
	\:.
\end{align}
This shows that we must compromise between resolution in time $\Delta t$ and resolution in frequency $\Delta f$, and how this is controlled by the segment length $T_0$. 
The balance between $\Delta f$ and $\Delta t$ is especially important when presenting data as spectrograms, as you will test out in the lab exercises.
 

\subsection{The Decibel-Scale}
A logarithmic scale allows visualization of a wider dynamic range than a linear scale and is preferred when the data set spans from very large to very small values. The decibel (dB) scale is the standard logarithmic scale in engineering and is defined as 
\begin{align}
	L_{dB}&= 10 \log_{10} \left(\frac{W}{W_{ref}} \right) = 20 \log_{10} \left(\frac{v}{v_{ref}} \right) 
\end{align}
where $v$ is an amplitude value (voltage, current, etc.) and $W$ is power or energy, so that $W/W_{ref}= (v/v_{ref})^2$. A value in \unit{dB} is always defined relative to a reference value, $v_{ref}$ or  $W_{ref}$. \qty{0}{dB} is this reference value. This reference can be a predefined value (e.g., \qty{1}{V} or \qty{1}{mW}), or an input or maximum value.

Spectral data are mostly presented in \unit{dB}. Since we often only interested in the relative variation between the spectral components, the reference value \qty{0}{dB} is often chosen as the maximum value in the data set. 
This results in all values in the spectrum or spectrogram being negative decibels, where \qty{0}{dB} is the maximum. The scale minimum gives the dynamic range, the span between the largest and smallest value presented. Typical values for this are \qty{-40}{dB} or \qty{-60}{dB}.

Decibels are never used for high-precision values.
Some important \unit{dB} values are listed in \tableref{tab:db-scale}.

\begin{table}[ht!]
	\caption{List of important \unit{dB}-values. }
	\label{tab:db-scale}
\begin{center}
\def\rspace{0.5em}
\begin{tabular}{cll}
	 	&	Amplitude ratio	& Power ratio \\ [\rspace]
	dB	&	$v/v_{ref}$	& $W/W_{ref}$	\\ [\rspace]
	\hline
	0	&	1				&	1			\\ [\rspace]
	-3	&	$\frac{1}{\sqrt{2}}= 0.71$	&	$\frac12 = 0.5$	\\  [\rspace]
	-6	&	$\frac12 = 0.5$			&	$\frac14 = 0.25$	\\	 [\rspace]
	-10	&	$\frac{1}{\sqrt{10}} = 0.32 \approx \frac{1}{3}$	&	$\frac{1}{10} = 0.1$	\\  [\rspace]
	-20	&	$\frac{1}{10} = 10^{-1} $			&	$\frac{1}{100} = 10^{-2}$	\\  [\rspace]
	-40	&	$\frac{1}{100} = 10^{-2} $			&	$\frac{1}{10 000} = 10^{-4}$	\\  [\rspace]
	-60	&	$\frac{1}{1000} = 10^{-3} $			&	$\frac{1}{1 000 000} = 10^{-6}$	\\  [\rspace]
	\hline
	
\end{tabular}
\end{center}	
\end{table}

\subsection{Fourier Series of Square and Triangle Waves}

Two important signals are the \emph{square} and \emph{triangular} waves shown in \figref{fig:square_triangle}. The Fourier coefficients for these can be calculated from \eqref{eq:fourier-coefficients}, the results are 
\begin{subequations}
\begin{align}
	& \text{Square wave } x_s(t)& a_k&= \begin{cases}
		\frac{2}{j \pi k }	\qquad & k = \pm 1, \pm 3, \pm 5, \ldots \\
		0    					& k = 0, \pm 2, \pm 4, \pm 6, \ldots
	\end{cases} 
	\label{eq:square_wave_coefficients} \\
	& \text{Triangle wave }  x_t(t)& a_k&= \begin{cases}
	\frac{4}{j \pi^2 k^2 }(-1)^{(k-1)/2}	& k = \pm 1, \pm 3, \pm 5, \ldots \\
	0    					& k = 0, \pm 2, \pm 4, \pm 6, \ldots
	\label{eq:triangle_wave_coefficients} 
\end{cases}
\end{align}
\end{subequations}
The even components ($k=\pm2, \pm4 , \ldots$) vanish for both these wave. This comes out of  equation \eqref{eq:fourier-coefficients}, but can also be seen from symmetry. 

\begin{center}
\begin{figure}[!h]
	\pythonfig{../figs/square_triangle.pdf}
	\caption{Square and triangle waves with the waves from the 1st and 3rd Fourier coefficients added. The triangle wave is  better reproduced by a few Fourier coefficients than the square wave, and its Fourier coefficients decrease more rapidly as $k$ increases.}
	\label{fig:square_triangle}
\end{figure}
\end{center}


\section{Lab Exercises}

\subsection*{Reporting}
Collect answers and code in a \jupyterlab notebook. Export this to pdf and upload it to Canvas.

You may prefer to do some of the coding in another development tool, such as \spyder. This is actually  recommend, as the testing and debugging options are better in \spyder than in \jupyterlab. You can load the .py-files from from \spyder into \jupyterlab and export them as a separate pdf-file.

\subsection{Spectrum and Spectrogram of Sinusoids}
This exercise will find the spectra and spectrograms for the two signals generated in Lab 1, use the functions you wrote in this lab. 
The frequency contents of these two signals does not change with time, so the spectrogram only should contain horizontal lines.
The two signals are

\begin{enumerate}[a)]
	\item Single-frequency signal with amplitude $A$, frequency $f$, phase $\phi$ and duration  given as
	\begin{align*}
		A&=\num{e4} & f&= \qty{1.5}{MHz} &	\phi&=-\ang{45} & \text{Duration \qty{4e-6}{s}}
	\end{align*}
	Set the sample rate to exactly 32 samples per period.
	
	\item A signal which is the sum of three sinusoids described by \\	
	\def\rspace{0.1em}
	\begin{tabular}{lcccc}
		\\
							& $k$			& 1	& 	2	& 	3	\\ [\rspace] \hline
		Frequency			& $f_k$ [Hz]	& 0	&	100	& 250	\\ [\rspace]
		Complex amplitude	&  $X_k$		& $10$ 	& 	$14e^{-j\pi/3}$		&	$8j$	\\ [\rspace] 
		\hline \\
	\end{tabular} \\
	Set the sample rate to \qty{10 000}{Samples/s} and duration to \qty{0.1}{s}.
\end{enumerate}

\subsubsection*{Exercises}	
\begin{enumerate}[1)]
	\item Plot the two signals and their Fourier coefficients $a_k$. Use the code in \tableref{tab:calculate-fft-code} as template.
	
	\item Plot the spectrogram of the two signals, using the code in  \tableref{tab:spectrogram-code} as template. Select a segment length and maximum frequency that fits to the signals, you may try different vales.
	
\end{enumerate}


\subsection{Beat}

\begin{enumerate}[1)]

\item Write a Python function to generate a beat signal defined by
	\begin{align}
		x(t) = A_1 \cos (2\pi(f_c-f_{\Delta})t ) + A_2 \cos (2\pi(f_c+f_{\Delta})t ) \:.
	\end{align}
	Specify the signal by the two amplitudes $A_1$ and $A_2$, the centre frequency $f_c$ and the difference frequency $f_{\Delta}$, the sample rate and the signal duration.
		
	A template for the function header is given in \tableref{tab:beat-code}.		
	You can make this simple by calling the function \verb|make_summed_cos| from Lab 1.
	
	\input{../figs/table-beat-code}	
		
\item Test the function for the input values
	\begin{align*}
		A_1 &= 10 & &, & A_2=&10 & &, & f_c&=\qty{400}{Hz} & &, & f_{\Delta}&=\qty{10}{Hz}
			& &, & f_s&=\qty{11025}{Hz} & &.
	\end{align*}
	Set the duration of the signal to \qty{2.0}{s}, but plot only the first \qty{0.5}{s}.
	
	Plot the signal and its power spectrum in two subplots. 
	Scale the frequency axis so that the frequencies in the beat are clearly identified.
	
	Comment the result.

\item Play the beat signal as a sound by using the module \verb|sounddevice|. 
	Listen to the signal and comment how it sounds.

\item Change the difference frequency to $f_{\Delta}=\qty{5}{Hz}$ and \qty{2}{Hz}.
	Plot these signals and listen to the sound of them.

	Comment how this changed the signal in the time and frequency domains. 

\item Use your function to generate a new beat signal with the following values
\begin{align*}
	A_1 &=  A_2=10 & &, & f_c&=\qty{1000}{Hz} & &, & \Delta f&=\qty{32}{Hz}
	& &, & f_s&=\qty{11025}{Hz} & &.
\end{align*}
Set the duration of the signal to \qty{0.26}{s} and plot the signal and its spectrum.

\item Plot the spectrogram of this signal using the template from \tableref{tab:spectrogram-code}.
Set the segment length to 1024 samples and limit the frequency scale so that the spectral lines are clearly seen. 

Comment the result.	

\item Do the same for segment lengths 512, 256, and 128 samples.

Comment how the different segment lengths change the appearance of the spectrogram. 
Relate this to the resolution in time and frequency from Section \ref{tf-resolution}.

What is the minimum segment length that can separate the two frequency lines?

\end{enumerate}


\subsection{Chirp}
This exercise repeats some of the tasks from the previous exercise with a chirp instead of a beat.
\begin{enumerate}
\item Write a Python function to generate a chirp signal defined by
		\begin{align}
			x(t)&= \cos \Psi(t) & 
			\Psi(t) &= 2 \pi \mu t^2 + 2 \pi f_0 t + \phi  \:.
		\end{align}
	Specify the signal by start and end frequencies, $f_1$  and $f_2$, phase $\phi$, sample rate, and duration. 
	
	A template for the header of the function is given in \tableref{tab:chirp-template-code}.
	
	\input{../figs/table-chirp-template-code.tex}
		
\item Test the function for the input values
	\begin{align*}
		f_1&=\qty{5000}{Hz} & &, & f_2&=\qty{300}{Hz}
		& &, & f_s&=\qty{11025}{Hz} & &.
	\end{align*}
	Set the duration of the signal to \qty{3.0}{s}.
	
	Generate the chirp signal and play it using the module \verb|sounddevice|. 	
	Comment how the signal sounds compared to the specification of the chirp.
		
\item Plot the spectrogram of this signal using the template from \tableref{tab:spectrogram-code}.
	Set the segment length to \num{2048} samples.
	Comment the result. Is it as expected?	
	
\item Do the same as above for segment lengths 1024, 512, 256, 128, 64, and 4096 samples.
	
	Comment how the different segment lengths change the appearance of the spectrogram.
	
	What seems to be the best segment length to resolve this chirp in time and frequency?

\item  Generate a new chirp with duration \qty{4}{s} starting at $f_1$=\qty{100}{Hz} and ending at $f_2$=\qty{4000}{Hz}. Set the sample rate to \qty{5000}{Hz}.

	Play the sound of this chirp using \verb|sounddevice| and plot the spectrogram. Use segment length 512 points.
	
	Comment the result.
	
	Change the sample rate to \qty{10000}{Hz} and repeat the tasks above. Since the sample rate is doubled, the segment length should also be doubled to get segments with the same duration in time.
	
	Comment the result. Why do the spectrograms look different?
	

	
\item Generate a new chirp with duration \qty{3}{s} starting at $f_1$=\qty{3000}{Hz} and ending at the negative frequency $f_2$=\qty{-2000}{Hz}.
	
	Listen to the signal. How does the frequency of the sound change?
	
	Plot the spectrogram of this chirp and explain the result.
	
	It may be easier to interpret this result if it is displayed as a \emph{two-sided} spectrogram that shows both positive and negative frequencies.	
	This is done by setting the parameter \verb|return_onesided| to \verb|False| and then arrange the negative and positive frequencies correctly using \verb|fftshift|.
	The code for this is shown below.	
	
\begin{lstlisting}[style=pythonstyle]		
f, t, sx = spectrogram(x, fs, nperseg=n_segment, return_onesided=False)
f = fftshift(f)
sx = fftshift(sx, axes=0)
\end{lstlisting}
	
\end{enumerate}

\subsection{Spectra of Square and Triangle Waves}	
\begin{enumerate}[1)]
\item Generate a square wave and a triangle wave, both with frequency \qty{100}{Hz}. 
	Each wave shall have exactly 2 periods. Use at least 100 samples per period to plot the signals.
	
	Create the signals by using square and triangle wave functions in \scipy, \verb|scipy.signal.square| and \verb|scipy.signal.sawtooth|. Look in the documentation for \scipy to see how to configure them.
	
\item Calculate and plot the Fourier coefficients of the two waves, use the template code in \tableref{tab:calculate-fft-code}.  

Compare the result with the values in \eqref{eq:square_wave_coefficients} and \eqref{eq:triangle_wave_coefficients}.

\item Calculate and plot the power spectrum of the two waves, use the template code in \tableref{tab:psd-code}.
Plot the spectrum on a decibel-scale.

Compare this result with the previous result.

\item Calculate and plot the spectrograms of the two waves, use the template code in \tableref{tab:spectrogram-code}.
Plot the intensity on a decibel-scale.

Compare this result with the previous results and comment.

\item Generate a chirp from a square wave. This is done by replacing the cosine-function in the chirp with the square-wave function. 

Let the frequency of the chirp start at \qty{100}{Hz} and end at \qty{4000}{Hz}. Set the sample rate to \qty{11025}{Hz}, and plot the spectrogram on a dB-scale.
A two-sided spectrogram may make the interpretation easier.

Play the sound of the chirp.

Comment the result. 

\end{enumerate}
	

\bibliographystyle{ieeetr}
\bibliography{../defs/tse2280.bib}
\end{document}